\section{Results}
\label{sec:results}

\begin{figure*}[t]
\centering
\setlength{\unitlength}{0.62mm}
\begin{picture}(270,50)
  \put(0,44){\textbf{(a)} Same-target sources}
  \put(5,25){\framebox(94,13){Fresh source: $q=t,\ k=\FreshOffset$}}
  \put(5,4){\framebox(94,13){Old\OldDelay{} source: $q=t-\OldDelay,\ k=\OldDelay$}}
  \put(119,29){\makebox(0,0){$q+k=t$}}
  \put(119,21){\circle*{2}}
  \put(99,31){\vector(1,-1){16}}
  \put(99,10){\vector(1,1){16}}
  \put(130,31){\framebox(63,12){arm selects source}}
  \put(130,2){\framebox(63,12){gripper selects source}}
  \put(119,21){\vector(1,1){11}}
  \put(119,21){\vector(1,-1){11}}
  \put(193,37){\vector(3,-1){45}}
  \put(193,8){\vector(3,1){45}}
  \put(238,22){\circle*{3}}
  \put(238,29){\makebox(0,0){physical target $t$}}
  \put(161.5,46){\makebox(0,0){each group independently selects an eligible same-target source}}
\end{picture}

\begin{minipage}[t]{0.59\textwidth}
\centering
\setlength{\unitlength}{0.60mm}
\begin{picture}(170,113)
  \put(0,106){\textbf{(b)} Frozen $2\times2$ factorial}
  \put(68,96){\makebox(0,0){Fresh gripper}}
  \put(108,96){\makebox(0,0){Old\OldDelay{} gripper}}
  \put(88,84){\makebox(0,0){\shortstack{$\textsc{FO20}-\textsc{Fresh}:\ \PosthocGripFreshArmDelta$ pp\\paired CI $\PosthocGripFreshArmPairedCI$}}}

  \put(50,55){\framebox(36,17){\shortstack{\textsc{Fresh}\\\ConfirmFreshRate\%}}}
  \put(90,55){\framebox(36,17){\shortstack{\textsc{FO20}\\\ConfirmFORate\%}}}
  \put(50,28){\framebox(36,17){\shortstack{\textsc{Reverse20}\\\ConfirmReverseRate\%}}}
  \put(90,28){\framebox(36,17){\shortstack{\textsc{FullOld20}\\\ConfirmFullOldRate\%}}}
  \put(46,63){\makebox(0,0)[r]{Fresh arm}}
  \put(46,36){\makebox(0,0)[r]{Old\OldDelay{} arm}}

  \put(46,51){\makebox(0,0)[r]{\shortstack{$\textsc{Reverse20}-\textsc{Fresh}$\\\PosthocArmFreshGripDelta{} pp\\paired CI $\PosthocArmFreshGripPairedCI$}}}
  \put(129,51){\makebox(0,0)[l]{\shortstack{$\textsc{FullOld20}-\textsc{FO20}$\\\PosthocArmOldGripDelta{} pp\\paired CI $\PosthocArmOldGripPairedCI$}}}
  \put(88,18){\makebox(0,0){\shortstack{$\textsc{FullOld20}-\textsc{Reverse20}:\ \PosthocGripOldArmDelta$ pp\\paired CI $\PosthocGripOldArmPairedCI$}}}
  \put(88,2){\makebox(0,0){post-hoc interaction $=\PosthocInteractionDelta$ pp, CI $\PosthocInteractionPairedCI$}}

  \put(137,78){\framebox(30,14){\shortstack{\textsc{hard-h\HardHorizon}\\\ConfirmHardRate\%}}}
  \put(152,97){\makebox(0,0){detached reference}}
\end{picture}
\end{minipage}\hfill
\begin{minipage}[t]{0.39\textwidth}
\centering
\setlength{\unitlength}{0.58mm}
\begin{picture}(112,113)(-36,0)
  \put(-36,106){\textbf{(c)} Cross-cohort paired effects}

  \put(-35,91){\makebox(0,0)[l]{Object development}}
  \put(-2,77){\makebox(0,0)[r]{C2$-$Fresh}}
  \put(-2,64){\makebox(0,0)[r]{h\HardHorizon$-$C2}}
  \put(-35,51){\makebox(0,0)[l]{Frozen confirmation}}
  \put(-2,37){\makebox(0,0)[r]{C2$-$Fresh}}
  \put(-2,24){\makebox(0,0)[r]{h\HardHorizon$-$C2}}

  \put(0,14){\line(1,0){75}}
  \put(25,18){\line(0,1){70}}
  \multiput(0,12)(25,0){4}{\line(0,1){4}}
  \put(0,7){\makebox(0,0){$-25$}}
  \put(25,7){\makebox(0,0){$0$}}
  \put(50,7){\makebox(0,0){$+25$}}
  \put(75,7){\makebox(0,0){$+50$}}
  \put(37.5,0){\makebox(0,0){paired effect (pp)}}

  \put(\FigObjCtwoLow,77){\line(1,0){\FigObjCtwoWidth}}
  \put(\FigObjCtwoLow,75){\line(0,1){4}}
  \put(\FigObjCtwoHigh,75){\line(0,1){4}}
  \put(\FigObjCtwoPoint,77){\circle*{3}}

  \put(\FigObjHardCtwoLow,64){\line(1,0){\FigObjHardCtwoWidth}}
  \put(\FigObjHardCtwoLow,62){\line(0,1){4}}
  \put(\FigObjHardCtwoHigh,62){\line(0,1){4}}
  \put(\FigObjHardCtwoPoint,64){\framebox(3,3){}}

  \put(\FigConfCtwoLow,37){\line(1,0){\FigConfCtwoWidth}}
  \put(\FigConfCtwoLow,35){\line(0,1){4}}
  \put(\FigConfCtwoHigh,35){\line(0,1){4}}
  \put(\FigConfCtwoPoint,37){\circle*{3}}

  \put(\FigConfHardCtwoLow,24){\line(1,0){\FigConfHardCtwoWidth}}
  \put(\FigConfHardCtwoLow,22){\line(0,1){4}}
  \put(\FigConfHardCtwoHigh,22){\line(0,1){4}}
  \put(\FigConfHardCtwoPoint,24){\framebox(3,3){}}
\end{picture}
\end{minipage}
\caption{Same-target temporal assignments expose component-dependent structure.
(a) Fresh and old sources both predict the physical target $t$; source age and
chunk offset are coupled by $q+k=t$, while arm and gripper can select different
eligible sources. (b) Frozen confirmation success rates in the factorial and a
detached fixed-h16 reference. The surrounding differences and interaction are
post-hoc decompositions with paired 95\% bootstrap intervals. Task-cluster
intervals are reported in the text, especially when interval types disagree.
The assignments do not independently randomize observation age and chunk
offset. (c) Paired effects and paired 95\% bootstrap intervals for the Object
development and frozen confirmation cohorts. Cohorts are not pooled. They use
different policy/checkpoint settings and are shown as separate paired effects;
the Object \textsc{hard-h16}$-$C2 completeness inference is post-hoc. Circles
denote C2$-$Fresh and squares denote \textsc{hard-h16}$-$C2.}
\label{fig:main}
\end{figure*}


\subsection{Q1: Does the Same Staleness Affect Components Differently?}

Moving the same one-second-old prediction from the gripper to the arm reduced
success from 83/140 (59.29\%) to 38/140 (27.14\%), a 32.14-point difference
(Fig.~\ref{fig:primary}). Under the sign convention
$S(\mathrm{A0G20})-S(\mathrm{A20G0})$, discordance was $48{:}3$ and the exact
McNemar test gave $p=1.96749\times10^{-11}$. The paired 95\% interval was
$[+23.6,+40.7]$ points, and the task-cluster interval was
$[+21.4,+44.3]$ points. This preregistered diagonal contrast establishes that
assigning the stale relationship to the arm or gripper is not behaviorally
interchangeable, but heterogeneous conditional effects prevent a unique
additive arm/gripper attribution.

The same asymmetry appeared when we later completed the previously missing
Spatial comparison. Fresh-arm/stale-gripper execution succeeded on 40/100
blocks and stale-arm/fresh-gripper execution on 12/100, a $+28.00$-point
difference with paired interval $[+17.00,+39.00]$ and task-cluster interval
$[+13.00,+45.00]$. This separate comparison preserves the sign but is not
pooled into the original 140-block result.

% Draft content placeholder only; final paper-facing artwork is intentionally absent.
\begin{figure*}[t]
\centering
\fbox{\parbox[c][31mm][c]{0.96\textwidth}{\small
\textbf{Draft Fig. 2 specification.} Left: absolute Object-126 success curves
for stale arm/fresh gripper and fresh arm/stale gripper at 0.10, 0.20, 0.40,
0.60, 0.80, 1.00, and 1.60~s. Mark 1.00~s as the originally frozen lag, not
an optimum. Right: five block-matched 1.00-s values, with distinct encodings
for historical anchors and later probes: translation stale 8.73\%, full arm
stale 9.52\%, rotation stale 42.06\%, Fresh 44.44\%, and gripper stale
64.29\%. Annotate translation-minus-rotation as $-33.33$ points.}}
\caption{\textbf{Age dependence and within-arm component structure on the
Object development cohort.} The fresh-arm/stale-gripper branch remains above
the stale-arm/fresh-gripper branch at every tested age, with observed
separations from $+21.43$ to $+54.76$ points; both curves are non-monotone.
The 1.00-s separation is the largest observed on the grid, but no
lag-versus-lag inference establishes a distinct peak. Historical anchors and
reviewer-directed probes are exactly block matched but were generated at
different study stages. The right panel is therefore a comparative
characterization, not a single five-condition preregistered factorial. The
stale-gripper effect magnitude should not be generalized across cohorts.}
\label{fig:development}
\end{figure*}


\subsection{Q2: How Does Sensitivity Vary with Age and Component?}

The fresh-arm/stale-gripper branch remained above the stale-arm/fresh-gripper
branch at every tested age from 0.10 to 1.60~s on the Object development
cohort (Fig.~\ref{fig:development}, left). Observed separations ranged from
$+21.43$ to $+54.76$ points, with no ordering reversal and the largest values
between 0.60 and 1.00~s. At 1.00~s, the separation was the largest observed on
the grid, but no lag-versus-lag inference establishes a distinct peak.

Neither absolute branch was monotone. Stale-arm/fresh-gripper success was
38.89, 38.89, 28.57, 16.67, 18.25, 9.52, and 25.40\% as age increased through
0.10, 0.20, 0.40, 0.60, 0.80, 1.00, and 1.60~s. It reached its observed
minimum at 1.00~s and recovered at 1.60~s. Fresh-arm/stale-gripper success was
60.32, 65.08, 69.84, 69.84, 71.43, 64.29, and 62.70\%. It was highest at
0.80~s and lower at 1.00 and 1.60~s. These shapes do not define a tolerance
threshold for either component.

Translation staleness was substantially more damaging than rotation staleness
on the same development cohort (Fig.~\ref{fig:development}, right).
Translation-stale execution succeeded on 11/126 blocks (8.73\%), whereas
rotation-stale execution succeeded on 53/126 (42.06\%). The primary contrast
$S(\mathrm{translation\ stale})-S(\mathrm{rotation\ stale})$ was $-33.33$
points, with discordance $3{:}45$, exact $p=1.31259\times10^{-10}$, paired
interval $[-42.06,-24.60]$, and task-cluster interval
$[-44.44,-22.22]$ points. All frozen leave-one-task-out estimates were
negative. This result supports component-dependent temporal sensitivity within
the evaluated ACT policy on this cohort.

Rotation staleness incurred little detectable cost relative to Fresh on this
cohort. Rotation-stale success was 42.06\%, compared with 44.44\% for Fresh,
a $-2.38$-point contrast ($p=0.607$); the paired interval
$[-8.73,+3.97]$ and task-cluster interval $[-6.35,+1.59]$ both included zero.
The exactly matched 1.00-s characterization was 8.73\% for translation stale,
9.52\% for the full arm stale, 42.06\% for rotation stale, 44.44\% for Fresh,
and 64.29\% for gripper stale. Staling translation alone nearly reproduces the
degradation observed when the full arm is stale. Historical anchors and later
component probes were generated at different study stages, and the magnitude
of the stale-gripper simple effect varies across cohorts.

% Draft content placeholder only; final paper-facing artwork is intentionally absent.
\begin{figure}[t]
\centering
\fbox{\parbox[c][24mm][c]{0.92\columnwidth}{\small
\textbf{Draft Fig. 3 specification.} Focal ordering contrast:
same-target dispersion $\mathrm{rotation}>\mathrm{translation}$, but
behavioral sensitivity $\mathrm{translation}>\mathrm{rotation}$. Secondary
rows: gripper persistence 0.675 at 1.00~s; occupancy $\rho=0.192$,
$p=0.309$; forecastability descriptive; command discontinuity omitted from
the main panel.}}
\caption{\textbf{Mechanism accounting.} Same-target disagreement predicts the
within-arm ordering incorrectly. Persistence, occupancy, forecastability, and
command-discontinuity diagnostics do not identify a positive causal mechanism.}
\label{fig:diagnostics}
\end{figure}


\subsection{Q3: Do Simple Diagnostics Explain the Ordering?}

Same-target prediction dispersion ranks the arm components in the opposite
order from behavioral sensitivity (Fig.~\ref{fig:diagnostics}). Normalized ACT
dispersion ranked rotation (0.1484) above translation (0.1364), whereas the
behavioral intervention ranked translation staleness as much more damaging
than rotation staleness. The disagreement measurement remains valid, but it
fails as an ordering predictor within the arm.

The remaining diagnostics did not recover a positive explanation. In training
demonstrations, the probability of no gripper transition within 1.00~s was
0.675, so persistence alone is insufficient. The prespecified gripper-activity
occupancy moderator was not supported ($n=30$, Spearman $\rho=0.192$,
$p=0.309$). Command-discontinuity comparisons cannot distinguish competing
explanations because conditions differ in their visited states and
trajectories. Forecastability curves were reported descriptively because no
curve pattern had been specified in advance as discriminative support, and
they did not provide a discriminative explanation. The tested diagnostics
therefore do not identify a positive causal mechanism.

\subsection{Q4: Does Component Structure Matter for Executable Schedules?}

Under dense matched querying, extending gripper commitment produced a clear
gain when the arm was also committed, whereas the corresponding change with a
fresh arm was small and uncertain. The four success rates were C00 (fresh arm,
fresh gripper) 77/140 (55.00\%), C10 (H16 arm, fresh gripper) 76/140
(54.29\%), C01 (fresh arm, H16 gripper) 81/140 (57.86\%), and C11 (H16 arm,
H16 gripper) 93/140 (66.43\%). With the arm committed, C11$-$C10 was
$+12.14$ points, with discordance $20{:}3$, exact $p=0.000488281$, paired
interval $[+5.71,+18.57]$, and task-cluster interval $[+3.57,+21.43]$.
With a fresh arm, C01$-$C00 was $+2.86$ points, with discordance $10{:}6$,
$p=0.454498$, paired interval $[-2.86,+8.57]$, and task-cluster interval
$[-4.29,+9.29]$.

Query schedule therefore cannot by itself explain this conditional effect in
the tested deterministic ACT evaluator. All four conditions queried the policy
at every executed step and discarded unused predictions. In a preregistered
identity check, dense C11 and sparse H16 produced exactly the same executed
actions, episode length, completion step, and terminal outcome from the same
initial state. This identity statement is limited to the tested deterministic
evaluator.

% Draft content placeholder only; final paper-facing artwork is intentionally absent.
\begin{figure*}[t]
\centering
\fbox{\parbox[c][30mm][c]{0.96\textwidth}{\small
\textbf{Draft Fig. 4 specification.} Main: pooled success versus log policy
queries per executed environment step for H16, H4, ARM4\_GRIP32, H2,
ARM2\_GRIP16, and TE\_DENSE. Use points only, no regression or connecting
line; give TE\_DENSE a distinct marker. Annotate H4$\rightarrow$ARM4\_GRIP32
($+4.667$ points) and H2$\rightarrow$ARM2\_GRIP16 ($+5.778$ points).
Identify H16 as strongest overall at 79.33\%. Inset: both paired effects for
LIBERO-10, Goal, and Spatial.}}
\caption{\textbf{Execution consequence at constrained replanning cadence.}
Extending gripper commitment improves success at two fixed arm cadences and
nearly matched policy-query rates per executed step. Coherent H16 remains the
strongest overall frozen operating point. The paired gains concentrate in
LIBERO-10; suite identity, baseline difficulty, and task semantics covary, so
the source of heterogeneity is not identified. Policy-query rate is not equal
compute.}
\label{fig:tracka}
\end{figure*}


Component-resolved commitment also improved executable schedules at two fixed
arm cadences (Fig.~\ref{fig:tracka}). Extending only gripper commitment from H4
to ARM4\_GRIP32 increased success from 314/450 to 335/450, or $+4.667$ points
(paired interval $[+2.00,+7.56]$; task-cluster interval
$[+0.67,+9.11]$). At the faster cadence, H2 to ARM2\_GRIP16 increased success
from 295/450 to 321/450, or $+5.778$ points (paired interval
$[+3.56,+8.22]$; task-cluster interval $[+2.67,+9.56]$). Policy-query rates
per executed step were nearly identical within each pair, approximately 0.251
and 0.501, respectively.

Coherent H16 nevertheless remained the strongest overall frozen operating
point at 357/450 successes (79.33\%). Component-resolved schedules therefore
have an operational consequence under constrained replanning cadence but are
not globally superior executors. Their gains were concentrated in LIBERO-10:
the H4 pair improved by 13.333 points there, 0.000 on Goal, and 0.667 on
Spatial; the H2 pair improved by 14.000, 2.000, and 1.333 points. H16 success
was 54.7\% on LIBERO-10, 91.3\% on Goal, and 92.0\% on Spatial. Suite identity,
baseline difficulty, and task semantics covary, so these data do not identify
the source of the concentration.

Canonical LeRobot v0.4.4 temporal ensembling is a secondary negative result.
With coefficient 0.01 and chunk length 100, TE\_DENSE heavily weighted old
predictions in this setup: its realized weighted mean age was 2.249~s, and it
produced substantially more near-boundary gripper commands than the other
frozen executors. This configuration is not an implementation bug, and the
near-boundary commands do not causally explain its full performance loss.
